\documentclass[11pt]{article}
\usepackage[T1]{fontenc}
\usepackage{lmodern,microtype}
\usepackage[margin=1in]{geometry}
\usepackage{amsmath,amsthm,amssymb,mathtools}
\usepackage{booktabs,array,enumitem}
\usepackage[hidelinks]{hyperref}
\hypersetup{pdftitle={Isolated prime clusters in every large modulus},
 pdfauthor={Meiji},pdfsubject={Least primes, arithmetic progressions, signed sieve weights}}
\newtheorem{theorem}{Theorem}[section]
\newtheorem{lemma}[theorem]{Lemma}
\newtheorem{proposition}[theorem]{Proposition}
\newtheorem{corollary}[theorem]{Corollary}
\theoremstyle{remark}
\newtheorem{remark}[theorem]{Remark}
\numberwithin{equation}{section}
\newcommand{\E}{\mathbb E}
\newcommand{\PP}{\mathbb P}
\newcommand{\Sg}{\mathfrak S}
\newcommand{\ind}{\mathbf 1}
\newcommand{\supp}{\operatorname{supp}}
\newcommand{\norm}[1]{\left\lVert#1\right\rVert}
\setlength{\emergencystretch}{2em}
\setlist[enumerate]{itemsep=3pt,topsep=5pt}
\title{\textbf{Isolated prime clusters in every large modulus}\\[5pt]
 \large Full support of first-prime locations and large gaps after a least prime}
\author{Meiji}
\date{25 November 2026\\[3pt]\small Research manuscript}
\begin{document}
\maketitle
\begin{abstract}
Put $T_Q=\phi(Q)\log Q$. For any fixed $0\le s<t<R$ and integer
$m\ge1$, we prove that a positive proportion of the reduced classes
modulo every sufficiently large integer $Q$ have at least $m$ primes
in $(sT_Q,tT_Q]$ and no primes elsewhere below $RT_Q$.
We also prove that a positive proportion have exactly one prime below
$RT_Q$, with that prime in the prescribed interval $(sT_Q,tT_Q]$.
Consequently every fixed-order prime has positive lower mass in every
nonempty interval at the first-prime scale. Even after restricting
the first prime to a fixed interval $sT_Q<p_1\le tT_Q$, with $s>0$,
the ratio $p_2/p_1$ is arbitrarily large on a positive proportion
of classes, while $p_m/p_1$ is arbitrarily close to one on another
positive proportion.
The proof extends the signed divisor-square method of Meiji's R119
manuscript. Its additional ingredient is a two-colour moment
factorization with an unmarked sieve majorant. This forces an empty
region and an occupied region in the same class while using at most
one prime indicator in each asymptotic calculation.
All moduli, including primorial moduli, are covered.
No prime-tuple hypothesis is assumed; exact occupancies above one,
a Poisson law, and a quantitative rate for growing parameters are
not asserted.
\end{abstract}

\section{Results and their relation to the least-prime problem}
\label{sec:results}

For an integer $Q\ge2$ and a reduced residue $a\bmod Q$, let
\[
 p_1(a,Q)<p_2(a,Q)<\cdots
\]
be the increasing sequence of primes congruent to $a$ modulo $Q$.
These sequences are infinite by Dirichlet's theorem. Write
\[
 N_Q(a,x)=\#\{p\le x:p\equiv a\pmod Q,\ p\text{ prime}\},
 \qquad T_Q=\phi(Q)\log Q.
\]
All probabilities give equal weight to the $\phi(Q)$ reduced classes.
The parameters in the following results are fixed before $Q$ tends
to infinity. Every positive lower-density constant and every threshold
for $Q$ may depend on those parameters, but never on the factorization
of $Q$.

\begin{theorem}[An isolated initial cluster]\label{thm:cluster}
Fix $0\le s<t<R$ and an integer $m\ge1$. There exist $d>0$ and $Q_0$ such
that, for every integer $Q\ge Q_0$, at least $d\phi(Q)$ reduced
classes satisfy
\begin{equation}\label{eq:cluster}
 \begin{split}
 N_Q(a,sT_Q)&=0,\\
 N_Q(a,tT_Q)&\ge m,\\
 N_Q(a,RT_Q)&=N_Q(a,tT_Q).
 \end{split}
\end{equation}
In particular their first $m$ primes all lie in $(sT_Q,tT_Q]$.
The number of primes in this interval is not prescribed exactly.
\end{theorem}

\begin{theorem}[A single isolated first prime]\label{thm:single}
Fix $0\le s<t<R$. There exist $d'>0$ and $Q_0'$ such that, for
every integer $Q\ge Q_0'$, at least $d'\phi(Q)$ reduced classes satisfy
\begin{equation}\label{eq:single}
 sT_Q<p_1(a,Q)\le tT_Q,\qquad p_2(a,Q)>RT_Q.
\end{equation}
Equivalently, exactly one prime in the class is at most $RT_Q$,
and it lies in the prescribed interval $(sT_Q,tT_Q]$.
\end{theorem}

For example, Theorem~\ref{thm:cluster} permits $m=10$,
$s=100$, $t=100.01$, and $R=1000$. It puts the first ten primes
of a positive proportion of classes into a window of width
$0.01T_Q$ near $100T_Q$ and excludes all primes in those classes
elsewhere below $1000T_Q$.
Theorem~\ref{thm:single} permits a first prime in
$(T_Q,1.01T_Q]$ and a second prime beyond $100T_Q$.
These examples concern different sets of classes and unspecified
positive proportions.

\subsection{Consequences at the natural scale}

\begin{corollary}[Positive local support]\label{cor:support}
For every fixed $j\ge1$ and $0\le s<t$,
\[
 \liminf_{Q\to\infty}
 \frac{\#\{a\bmod Q:(a,Q)=1,\ sT_Q<p_j(a,Q)\le tT_Q\}}{\phi(Q)}
 >0.
\]
More strongly, a positive proportion have all of their first $j$
primes in that interval.
\end{corollary}

\begin{corollary}[Ratios with a first prime at the natural scale]
\label{cor:ratios}
Fix $0<s<t$ and put
\[
 \mathcal R_Q(s,t)=\{a\bmod Q:(a,Q)=1,\
                         sT_Q<p_1(a,Q)\le tT_Q\}.
\]
For every fixed $m\ge2$,
\begin{equation}\label{eq:ratios}
 \min_{a\in\mathcal R_Q(s,t)}\frac{p_m(a,Q)}{p_1(a,Q)}
       \longrightarrow1,
 \qquad
 \max_{a\in\mathcal R_Q(s,t)}\frac{p_2(a,Q)}{p_1(a,Q)}
       \longrightarrow+\infty.
\end{equation}
Each fixed tolerance in the first limit and each fixed lower bound
in the second hold on a positive proportion of all reduced classes.
\end{corollary}

The restriction $p_1\asymp T_Q$ matters. A very small first prime
can make $p_2/p_1$ large for an elementary reason when $Q$ grows.
Corollary~\ref{cor:ratios} rules out that explanation and asserts
positive proportions. All consecutive primes and prime gaps in
this paper are consecutive \emph{within one reduced class}.

\subsection{The classical conjecture reached by the argument}

Let $P(Q)=\max_{(a,Q)=1}p_1(a,Q)$.
\begin{corollary}[All-modulus divergence]\label{cor:gp}
For every fixed $C>0$, a positive proportion of reduced classes
modulo every sufficiently large $Q$ satisfy $p_1(a,Q)>CT_Q$.
In particular
\begin{equation}\label{eq:gp}
 \frac{P(Q)}{\phi(Q)\log Q}\longrightarrow+\infty
 \qquad(Q\to\infty).
\end{equation}
\end{corollary}

Erd\H os \cite[p.~57]{E49} proposed the failure of an almost-all
upper bound for every fixed multiple of $T_Q$. His late-class
theorem gave one multiple on infinitely many moduli; his early-class
theorem supplied positive proportions below every fixed positive
multiple. The all-large-modulus formulation with one multiple
appears in \cite[p.~209]{E65} and is catalogued as problem 971
\cite{Bloom}.
Granville and Pomerance \cite[p.~194]{GP} explicitly separated
\eqref{eq:gp} from the stronger all-modulus conjecture
$P(Q)\gg\phi(Q)(\log Q)^2$. We obtain \eqref{eq:gp}, including
the difficult primorial sequence, but not that stronger bound.
The quantitative almost-all-modulus lower bound of Li, Pratt, and
Shakan \cite{LPS} addresses a different quantifier and is not
superseded by a numerical rate here.

Corollary~\ref{cor:gp} already follows from the main theorem of the
R119 working manuscript \cite{Meiji119}; it is not a new consequence
unique to this extension. The new step developed here is to couple
suppression and occupation in the \emph{same} class, and to isolate
exactly one prime at a prescribed location. All transfer estimates
and functional constructions needed for this extension are proved
below, so the mathematical argument does not depend on an omitted
proof in R119.

\subsection{Comparison with a Poisson prediction}

The prime number theorem gives, for every fixed $c>0$,
\begin{equation}\label{eq:firstmean}
 \E_a N_Q(a,cT_Q)
 =\frac{\pi(cT_Q)-\#\{p\mid Q:p\le cT_Q\}}{\phi(Q)}
 \longrightarrow c.
\end{equation}
Indeed $Q/\phi(Q)\ll\log\log Q$, so
$\log(cT_Q)/\log Q\to1$ and the excluded prime divisors contribute
$o(\phi(Q))$.
Thus $T_Q$ is the scale at which the \emph{mean count} is one.
This observation does not identify the mean of $p_1$.

Under a uniform prime-tuple hypothesis, Leung proves a Poisson
count law and an exponential law for the first prime
\cite{LeungMoments}. His point-process paper \cite{LeungProcess}
also obtains the joint Poisson probabilities for disjoint windows
under such a hypothesis. Our conclusions give unconditional positive
lower mass for particular events predicted by those laws. They do
not supply the predicted probabilities or independence between
several occupied windows. Averaging singular series over products
of intervals is already part of that literature; the bounded-range
version needed here is included with its uniformity in $Q$.

The signed functional construction is adapted from the ordinary-gap
method of \cite{OAI}, as in \cite{Meiji119}. The arithmetic input
is the classical Bombieri--Vinogradov theorem. We use at most one
prime-marked row in each asymptotic calculation; additional prospective primes are bounded by divisor-square
majorants, leaving at most one actual prime mark.

\paragraph{Proof structure.}
Sections~\ref{sec:mixed}--\ref{sec:forms} establish the arithmetic
transfer, coloured singular-series averages, and quadratic forms.
Section~\ref{sec:variational} gives arbitrarily strong suppression
and amplification with arbitrarily small fixed support budgets.
Section~\ref{sec:joint} proves the same-class coupling.
Section~\ref{sec:singleton} isolates one prime and passes from
rows to physical prime intervals.
Section~\ref{sec:consequences} proves the remaining consequences
and states the precise boundary of the results.


\section{An arithmetic mixed-moment transfer}\label{sec:mixed}

Fix $K>0$, let $L=\log Q$, $u=Q/\phi(Q)$, and
$H=\lfloor KL/u\rfloor$. All rows in this section belong to
$I_0=\{1,\ldots,H\}$. Enlarging $K$ by a fixed amount permits
any bounded $O(L/u)$ range of positive rows. We use the standard
bound $u\ll\log\log Q=\log L$, hence $L/u\to\infty$.
Throughout $a$ is uniform on the representatives $1\le a\le Q$ with
$(a,Q)=1$. Put $n_b=a+bQ$ and $\vartheta(n)=\ind_{n\;\mathrm{prime}}\log n$.
For fixed smooth tests $F_r$ in dimensions $j_r$, take compactly
supported smooth extensions from the nonnegative orthants to the full
spaces. On the orthant $F_r$ vanishes for $\sum x_i>\rho_r$. Define
\[
 D_{F_r}(a;b_1,\ldots,b_{j_r})=
 \sum_{d_i\mid n_{b_i}}\prod_i\mu(d_i)\,
 F_r((\log d_i/L)_i).
\]
A zero-dimensional test is a scalar. Across all tests prescribe a
partition of the coordinates $\Gamma$ into $t$ nonempty fibers
$M_1,\dots,M_t$, meaning coordinates in a fiber have the same row;
different fibers use distinct rows $b_1,\dots,b_t\in I_0$.
If $\delta=1$, allow one additional prime-marked row $b_0\in I_0$
distinct from these; if $\delta=0$ there is no mark. Write $J$ for
the set of $t+\delta$ row indices, $Y_0=1$, $Y_1=\vartheta(n_{b_0})$,
and
\[
 \Sg_Q(J)=\prod_{p\nmid Q}
 (1-\nu_p(J)/p)(1-1/p)^{-|J|},\qquad
 \nu_p(J)=\#(J\bmod p).
\]
\begin{lemma}\label{mixed}
Suppose $\rho=\sum_r\rho_r<1$ if $\delta=0$, and $\rho<1/2$ if
$\delta=1$. For fixed dimensions, tests and fiber partition, uniformly
over the distinct row indices,
\[
 \E\,Y_\delta\prod_r D_{F_r}
 =u^{t+\delta}L^{-t}
 \big\{\Sg_Q(J)\mathcal C+O(L^{-1/2}(\log L)^A)\big\}.
\]
Here $A$ is fixed and $\mathcal C$ depends only on the tests and
the partition. In particular no relative error is asserted when a local
obstruction makes $\Sg_Q(J)=0$.
\end{lemma}
\begin{proof}
Expand the divisor sums. All surviving tuples are squarefree in each
coordinate, have $\prod_{\gamma\in\Gamma}d_\gamma\le Q^\rho$,
and are coprime to $Q$. Write $D=\mathrm{lcm}(d_\gamma)$.
For each $p\mid D$, all its coordinates must have row index congruent
to one common class $b\bmod p$. This imposes one residue on $a$ modulo
$D$. Without a mark, its conditional density is
$D^{-1}+O(2^{\omega(Q)}/\phi(Q))$: expand $\ind_{(a,Q)=1}$ by
M\"obius inversion and apply CRT for each divisor of $Q$. The number of
tuples is $O(Q^\rho(1+L)^{|\Gamma|})$, by successive harmonic summation.
Since $2^{\omega(Q)}=Q^{o(1)}$, the total error is $O(Q^{\rho-1+o(1)})$.

With a mark, if a selected row $b$ agrees with $b_0\bmod p$, then
$p\mid n_{b_0}$ and the marked term is zero, since $n_{b_0}>Q>D$.
Otherwise $n_{b_0}$ is in a reduced class $r\bmod D$ within
$(b_0Q,(b_0+1)Q]$. A prime in this interval already ensures $(a,Q)=1$.
Thus the conditional marked mass is $u/\phi(D)$ plus an error.
For clarity, write $E_D$ for the maximum, over these intervals and
all reduced residues modulo $D$, of the absolute difference between the
sum of $\vartheta$ and $Q/\phi(D)$. For each fixed $\kappa_0\ge1,B>0$,
\[
 \sum_{D\le Q^\rho}\mu^2(D)\kappa_0^{\omega(D)}E_D\ll QL^{-B}.
\]
Indeed the maximal-endpoint Bombieri--Vinogradov theorem \cite{Bombieri,Vaughan} up to
$X=(H+1)Q\ll QL$ gives the unweighted estimate with any log saving,
since $Q^\rho\le X^{1/2}(\log X)^{-D_0}$ for every fixed $D_0$ eventually.
Take differences of two endpoint estimates. Replacing von Mangoldt
weights by prime-only weights costs at most
$O(\sqrt X(\log X)^2)$ per modulus, hence $O(Q^{1/2+\rho+o(1)})$ even
with weights. The trivial count in an interval gives
$E_D\ll Q L^2/D$ ($D\le Q^\rho<Q$, $D/\phi(D)\ll\log L$).
Therefore $\sum_{D\le Q^\rho}\mu^2(D)\kappa_0^{2\omega(D)}E_D
\ll Q L^{O_{\kappa_0}(1)}$, by the Euler product for the harmonic sum.
Cauchy--Schwarz between this bound and the unweighted estimate proves
the display. At most $2^{|\Gamma|\omega(D)}$ tuples have lcm $D$;
bounded tests and division by $\phi(Q)$ absorb their total error to
arbitrary negative powers of $L$. Only the moduli $D$, not $QD$,
enter this application of Bombieri--Vinogradov.

Write $w_\gamma=1+is_\gamma$, $z_\gamma=w_\gamma/L$. Smooth extension
and Fourier inversion give
$F_r(\mathbf x)=\int\Phi_r(\mathbf s)\exp(-\sum w_\gamma x_\gamma)d\mathbf s$
on the orthant, with Schwartz $\Phi_r$ (the Fourier normalization
is included in $\Phi_r$). The main local-density sum, before Fourier
inversion, was finite by support; one may equivalently extend it to
all tuples since the tests vanish outside that support.
After inversion the sum converges absolutely, uniformly in $\mathbf s$,
by $\prod_p(1+O_{|\Gamma|}(p^{-1-1/L}))\ll L^{O_{|\Gamma|}(1)}$.
There is an external factor $u^\delta$ and the local factors are
$P_p=1$ for $p\mid Q$ and otherwise
\[
 P_p=1+\frac1{p-\delta}
 \sum_{\varnothing\ne U\subseteq\Gamma}^{*}
 (-1)^{|U|}p^{-z_U},\quad z_U=\sum_{\gamma\in U}z_\gamma.
\]
The star requires all rows in $U$ to agree modulo $p$ and, if marked,
to avoid the marked row modulo $p$.

Separate the universal polar product via $H_p=P_p/B_p$, where
\[
 B_p=\prod_{i=1}^t\prod_{\varnothing\ne U\subseteq M_i}
 (1-p^{-1-z_U})^{-(-1)^{|U|}}.
\]
For $|s_\gamma|\le\sqrt L$, the pole expansion of $\zeta(1+z_U)$
gives
\[
 \prod_pB_p=L^{-t}\mathcal K(\mathbf s)(1+O(L^{-1/2})),\qquad
 \mathcal K=\prod_{i=1}^t\prod_{\varnothing\ne U\subseteq M_i}
 \left(\sum_{\gamma\in U}w_\gamma\right)^{(-1)^{|U|+1}}.
\]
Indeed the sum of $(-1)^{|U|}$ over nonempty subsets of one fiber is
$-1$. The kernel grows at most polynomially; the complementary Fourier
tail costs $O(L^{-B})$ for any fixed $B$ by the absolute Euler bound.

At $\mathbf z=0$, $B_p=(1-1/p)^t$. For each distinct allowed row
class the subset sum is $-1$; with a mark the number of allowed classes
is $\nu_p(J)-1$. It follows that
\[
 \prod_p H_p(0)=u^t\Sg_Q(J).
\]
We justify evaluation at zero even at locally obstructed tuples.
For $p\le L^2$, the factors $H_p(z)$ and $H_p(0)$ have modulus at most
$1+C/p$ (increase the fixed $C$ for small primes), and their
difference is $O(L^{-1/2}\log p/p)$ on the clipped Fourier region.
This follows by differentiating their finite local formulas along
the straight segment from $0$ to $z$; all denominators
$1-p^{-1-z_U}$ are bounded away from zero there. Telescoping the products
costs $O(L^{-1/2}(\log L)^A)$ by Mertens' estimates. For $p>L^2$
larger than every row difference, the first-order terms in $P_p$ and
$B_p$ agree if $p\nmid Q$, so $H_p=1+O(p^{-2})$. For $p\mid Q$ we only have
$1+O(1/p)$, but $\sum_{p\mid Q,\,p>L^2}1/p\le L/(L^2\log(L^2))$.
Both tails are negligible, also when multiplying the small-prime
products, which are at most $(\log L)^{O(1)}$. Thus the difference
from the zero-value product is bounded additively as above.
Integrating leaves
$\mathcal C=\int (\prod_r\Phi_r)\mathcal K\,d\mathbf s$.
Since $u\ge1$, normalization gives the asserted error.
\end{proof}


\section{Singular-series averages on coloured row sets}\label{sec:average}

Throughout the rest of the proof, a row set is a union of at most a
fixed number of integer intervals contained in
$[1,KL/u]$, where $K$ is fixed. We use disjoint sets $I_A,I_B$
with
\begin{equation}\label{eq:rowcounts}
 |I_i|=\lambda_i L/u+O(1),\qquad \lambda_i>0,\quad i=A,B.
\end{equation}
All implied constants and the number of interval components are fixed.
The endpoints may depend on $Q$. A single row set is also allowed.

\begin{lemma}[Coloured singular-series average]\label{av}
Fix nonnegative integers $\ell_A,\ell_B$. Choose ordered distinct
rows within each colour, with $\ell_i$ rows of colour $i$, and let
$J$ be their union. Then the mean of $\Sg_Q(J)$ is $1+o(1)$.
The error is uniform over the row sets just described.
\end{lemma}
\begin{proof}
The empty and singleton cases are exact. Write $\ell=\ell_A+\ell_B$,
and let $H_*$ be the smallest size of a colour set that is sampled.
Then $H_*\asymp L/u\to\infty$.
Put $y=(\log H_*)/4$ and
$P=\prod_{p\le y,\ p\nmid Q}p$. The prime number theorem gives
$P\le H_*^{1/2}$ for all sufficiently large $Q$.
A uniform sample from a union of at most $J_0$ integer intervals
differs modulo $P$ from uniform measure by total variation
$O(J_0P/H_*)$.

For independent uniform coordinates modulo $p$, introduce another
independent uniform residue and ask that it avoid all $\ell$
coordinates. This proves
\[
 \E(1-\nu_p(J)/p)=(1-1/p)^\ell.
\]
The Chinese remainder theorem therefore gives mean one for the
truncated singular series over independent uniform coordinates
modulo $P$. The truncated product is at most
$(\log y)^{O_\ell(1)}$ by Mertens' estimates.
Changing the sampling distribution costs
$O_\ell(P/H_*)$ times this bound. Removing repeated coordinates
within a colour costs $O_\ell(1/H_*)$ times the same bound.
Different colours cannot share a row because their row sets are
disjoint.

For any resulting tuple of distinct rows put
$\Delta=\prod_{i<j}|b_i-b_j|$. Then
$\log\Delta=O_\ell(\log L)$.
For $p>y>2\ell$, a tail factor has logarithm
$O_\ell(p^{-2})$ if $p\nmid\Delta$, and
$O_\ell(p^{-1})$ otherwise. Since each such prime divisor has
$\log p\ge\log y$,
\[
 \sum_{\substack{p>y\\p\mid\Delta}}\frac1p
 \le \frac{\log\Delta}{y\log y}.
\]
Here $\log H_*\sim\log L$, because $u\ll\log L$. Thus the total
tail logarithm is
$O_\ell(y^{-1}+\log\Delta/(y\log y))=o(1)$ uniformly.
Omitting primes dividing $Q$ only improves this estimate.
The tail is $1+o(1)$ multiplicatively, preserving any zero factor
in the truncated product. This proves the assertion.
\end{proof}

\begin{remark}\label{rem:factor}
In Lemma~\ref{mixed}, every fiber of coordinates belonging to
one row has a single colour. If every test depends only on rows
of one colour, the Fourier kernel $\mathcal K$ factors across the
colours. Its integral against the product Fourier transform then
factors as well. Lemma~\ref{av} removes the singular-series coupling
after the row labels are averaged. This gives factorization of the
specific sieve moments below; it is not an independence assertion
about the prime counts.
\end{remark}

\section{Quadratic weights and unmarked prime detectors}\label{sec:forms}

Fix a support budget $\tau>0$. For a row set $I$ with
$|I|u/L\to\lambda>0$, take a finite family consisting of a real
scalar $f_0$ and real symmetric smooth functions $f_j$,
$1\le j\le k$, compactly supported in
\[
 \{\mathbf t:t_i>0,\quad t_1+\cdots+t_j<\tau\}.
\]
Put $f_{k+1}=0$ and
\[
 F_j(\mathbf x)=\int_{\substack{t_i\ge x_i\\1\le i\le j}}
                     f_j(\mathbf t)\,d\mathbf t,
 \qquad
 (Tf_{j+1})(\mathbf t)=\int_0^\infty f_{j+1}(\mathbf t,z)\,dz,
 \qquad F_0=f_0 .
\]
The cumulative tests $F_j$ have smooth compact extensions to the
full spaces: extend $f_j$ by zero, integrate, and multiply by
smooth cutoffs on the negative half-lines. On the nonnegative
orthant their signed complete derivatives are
$(-1)^j\partial_1\cdots\partial_jF_j=f_j$.
It is these derivatives, rather than the cumulative tests, that
vanish on every coordinate face in positive dimensions.

Define
\begin{align}
 Z&=\sum_{j=0}^k\ \sum_{\substack{S\subset I\\|S|=j}}
          D_{F_j}(a;S),&
 V&=\frac1L\sum_{b\in I}\vartheta(a+bQ),&
 N&=\sum_{b\in I}\ind_{a+bQ\ {\rm prime}},\label{eq:ZVN}\\
 \alpha_j&=\frac{\lambda^j}{j!},&
 w&=\sum_{j=0}^k\alpha_j\norm{f_j}_2^2,&
 v&=\lambda\sum_{j=0}^k\alpha_j\norm{f_j+Tf_{j+1}}_2^2.
 \label{eq:wv}
\end{align}
Uniformly on the allowed positive rows,
\begin{equation}\label{eq:NV}
 N\le V\le(1+\epsilon_Q)N,\qquad
 \epsilon_Q=O_K(\log L/L)\longrightarrow0.
\end{equation}

\begin{lemma}[The quadratic forms]\label{forms}
If $2\tau<1/2$, then
\[
 \E Z^2\longrightarrow w,\qquad
 \E Z^2V\longrightarrow v.
\]
For each fixed even integer $r$ with $r\tau<1$,
$\E|Z|^r=O(1)$. For each fixed integer $r\ge1$,
$\E V^r=O_{r,\lambda,K}(1)$.
The constants for $Z$ may depend on the entire fixed family.
\end{lemma}
\begin{proof}
In a product of two subset tests, each fiber has size one or two.
The corresponding factor in $\mathcal K$ is $w_1$ for size one
and
\[
 \frac{w_1w_2}{w_1+w_2}
   =\int_0^\infty w_1w_2e^{-(w_1+w_2)t}\,dt
\]
for size two. Fourier inversion therefore expresses the kernel
constant as an integral of the two signed complete derivatives:
a shared coordinate is integrated over a common variable, while
an unpaired coordinate is evaluated at zero. The latter makes
the constant vanish. Equal subsets of size $j$ contribute
$\norm{f_j}_2^2$. All interchanges are absolute because the
Fourier transforms are Schwartz and $\Re w_i=1$.
After Lemmas~\ref{mixed} and \ref{av}, the diagonal contribution is
\[
 \binom{|I|}{j}(u/L)^j(1+o(1))\norm{f_j}_2^2
       \longrightarrow\alpha_j\norm{f_j}_2^2.
\]

For a marked row $b_0$, primality of $a+b_0Q>Q$ implies that
the only allowed divisor of that row in a supported test is one.
Split the subsets according to membership of $b_0$.
On the marked event the new dimension-$j$ test is
$F_j+F_{j+1}(\cdot,0)$, whose signed complete derivative is
$f_j+Tf_{j+1}$. This again vanishes on positive-dimensional
coordinate faces. There are
$|I|\binom{|I|-1}{j}$ diagonal labelled terms, and the marked
normalization from Lemma~\ref{mixed}, divided by $L$, gives
the limit $\lambda^{j+1}/j!$. This proves the formula for $v$.

For any fixed product of tests there are only finitely many
coincidence patterns. A pattern with $t$ distinct unmarked rows
has $O((L/u)^t)$ labelings. Its total error from
Lemma~\ref{mixed} is
\[
 O\bigl(L^{-1/2}(\log L)^A\bigr)=o(1).
\]
For a prime mark there is one additional row factor $L/u$
and one additional factor $u/L$, so the same estimate holds.
All fixed kernel integrals are finite. Nonnegative singular-series
sums are bounded by Lemma~\ref{av}, after dropping ordering
constraints if necessary. This proves the even-moment assertion
when the sum of the $r$ support budgets is less than one.

For the prime-count moment, choose a real smooth univariate test
$G_r$ with $G_r(0)=1$, supported on the orthant below a fixed
$\sigma_r$ satisfying $2r\sigma_r<1$.
For a prime row $D_{G_r}(a;b)=1$. Hence
\[
 N\le \sum_{b\in I}D_{G_r}(a;b)^2.
\]
The $r$-th power of the right side expands into $2r$ unmarked
tests, with total support budget less than one.
The same finite-pattern argument bounds its expectation.
Equation~\eqref{eq:NV} gives the bound for $V$.
\end{proof}

We now fix, for the remaining constructions,
\begin{equation}\label{eq:budgets}
 \tau=\sigma=1/64.
\end{equation}
Choose a real smooth test $G$ with $G(0)=1$ and
$G(x)=0$ for $x\ge\sigma$ on the orthant. Write
\begin{equation}\label{eq:detectors}
 \begin{split}
 \kappa&=\int_0^\infty |G'(x)|^2\,dx>0,\qquad
 U_B=\sum_{b\in I_B}D_G(a;b)^2,\\
 R_B&=\sum_{\substack{b,c\in I_B\\b\ne c}}
             D_G(a;b)^2D_G(a;c)^2 .
 \end{split}
\end{equation}
The last sum is over ordered distinct rows. Pointwise,
\begin{equation}\label{eq:majorants}
 N_B\le U_B,\qquad
 N_B(N_B-1)\le R_B,\qquad
 V_B\le(1+\epsilon_Q)U_B.
\end{equation}

\begin{lemma}[Two-colour moments]\label{lem:tensor}
Give each colour its own fixed family $Z_A,Z_B$ and parameters
$w_A,v_A,w_B,v_B$ as in \eqref{eq:wv}. There is a finite constant
$q_B\ge0$, depending only on the $B$ family, $\lambda_B$, and $G$,
such that, with $W=Z_A^2Z_B^2$,
\begin{align}
 \E[Z_B^2U_B]&\longrightarrow q_B,&
 \E W&\longrightarrow w_Aw_B,\label{eq:tensor1}\\
 \E[WV_B]&\longrightarrow w_Av_B,&
 \E[WV_AU_B]&\longrightarrow v_Aq_B.\label{eq:tensor2}
\end{align}
Moreover $\E W^2=O(1)$ and $\E V_B^4=O(1)$.
For a family only on $A$, the additional detector identities are
\begin{align}
 \E[Z_A^2N_B]&\longrightarrow w_A\lambda_B,\label{eq:singlefirst}\\
 \E[Z_A^2R_B]&\longrightarrow w_A\kappa^2\lambda_B^2,\label{eq:singlepair}\\
 \E[Z_A^2V_A]&\longrightarrow v_A.\label{eq:singlemark}
\end{align}
\end{lemma}
\begin{proof}
Expand each expression into its finite set of row-coincidence
patterns. By Remark~\ref{rem:factor}, the kernel constants factor
over the two colours, and Lemma~\ref{av} gives mean singular
series one for each prescribed colour pattern.
Two subset tests on an unmarked colour contribute its quadratic
constant $w_i$, by the orthogonality calculation in
Lemma~\ref{forms}. A marked colour contributes $v_i$, by the
same marked-row deletion argument.

In $Z_B^2U_B$, the detector row may coincide with a row of either
subset test. Each resulting pattern has a fixed finite kernel
integral and an asymptotic count of row labels. Their sum has a
finite limit $q_B$. It is nonnegative because $Z_B^2U_B\ge0$.
Only the $B$ family occurs in this definition.
For $WV_AU_B$, the sole prime mark has colour $A$, while the
$B$ factor is entirely unmarked. The same pattern sum therefore
factors into $v_Aq_B$. The other identities in
\eqref{eq:tensor1}--\eqref{eq:tensor2} follow identically.
Their worst marked support budget is
$2\tau+2\tau+2\sigma=6/64<1/2$.
The errors summed over all labels are $o(1)$ as in
Lemma~\ref{forms}.

For \eqref{eq:singlefirst}, the $B$ colour has just a distinct
prime mark and no subset test. Its contribution after summing
over $B$ is $\lambda_B$.
One first obtains this identity with $V_B$; replacing $V_B$
by $N_B$ costs $o(1)$ using \eqref{eq:NV} and the bounded
weighted first moment.
For \eqref{eq:singlepair}, the two distinct $B$ rows each carry
two copies of $G$. The two-row kernel is
\[
 \left(\int_0^\infty |G'(x)|^2\,dx\right)^2=\kappa^2.
\]
The ordered row count contributes $\lambda_B^2$.
This is an asymptotic for a divisor majorant, not for a pair
of prime indicators.
Equation~\eqref{eq:singlemark} is Lemma~\ref{forms}.

Finally $W^2=Z_A^4Z_B^4$ has total unmarked support budget
$8\tau=1/8<1$.
The bound for $V_B^4$ follows from $U_B^4$, whose eight tests
have total support budget $8\sigma=1/8<1$.
These bounds again follow from the finite-pattern argument.
\end{proof}


\section{Arbitrarily strong suppression and amplification}
\label{sec:variational}

This section is purely functional and independent of $Q$.
For each fixed $\lambda>0$ define $w,v$ by \eqref{eq:wv}.

\begin{lemma}[The two variational directions]\label{lem:variational}
For every support budget $\tau>0$ the following hold.
\begin{enumerate}[label=(\roman*)]
 \item For every $\eta>0$, there is a finite interior-supported
 family with $w>0$ and $v<\eta w$.
 \item For every $H_0>0$, there is a finite interior-supported
 family with $w>0$ and $v>H_0w$.
\end{enumerate}
The family is fixed before the arithmetic limit $Q\to\infty$.
\end{lemma}
\begin{proof}
We include the signed cancellation construction underlying
\cite{OAI,Meiji119}.
For any $M>0$, choose a nonnegative
$g\in C_c^\infty((0,\infty))$ with
\begin{equation}\label{eq:g}
 \int g^2=1,\qquad \int g=\sqrt M,\qquad
 \mu_g:=\int t g(t)^2\,dt<\beta,
 \qquad 0<\beta<\gamma<\tau.
\end{equation}
Here is an explicit existence argument.
Take a smooth nonnegative cutoff of $1/t$ equal to $1/t$ on
$[1,R_0]$, bounded by $1/t$, and supported in $[1/2,2R_0]$.
Normalize in $L^2$ and then dilate, preserving its $L^2$ norm,
to prescribe the integral $\sqrt M$.
If $I_1$ and $I_t$ are respectively the integral and the
$t$-weighted squared integral of the original cutoff, the
new first moment is $MI_t/I_1^2$.
It is at most
$M\log(4R_0)/(\log R_0)^2$, which tends to zero.
This proves \eqref{eq:g}.

Let $B_0$ exceed the support endpoint of $g$, and choose smooth
$0\le\chi\le1$, equal to one for $0\le x\le\beta$ and zero
for $x\ge\gamma$. For an integer $k$ and $0\le j\le k$, put
\[
 P_j(\mathbf t)=k^{j/2}\prod_{i=1}^j g(kt_i),
 \qquad S_j=t_1+\cdots+t_j,\qquad P_0=1.
\]
Then $\norm{P_j}_2=1$ and
$TP_{j+1}=\sqrt{M/k}\,P_j$.
Under the probability density $P_j^2$, the mean of $S_j$
is $j\mu_g/k\le\mu_g$ and its variance is at most
$\operatorname{Var}_{g^2}(t)/k$.
Chebyshev's inequality, with the fixed gap $\beta-\mu_g>0$,
gives
\begin{equation}\label{eq:cutoffexception}
 \PP_j\{S_j>\beta-B_0/k\}=O_g(k^{-1})
 \quad(0\le j\le k).
\end{equation}
Off this exceptional set all cutoffs remain one, even after
integrating over one extra coordinate in the support of $g$.

\paragraph{Suppression.}
Take $M=\lambda$, let $r=\lfloor\sqrt{k}\rfloor$, and set
\[
 f_j=(-1)^j(r\alpha_j)^{-1/2}P_j\chi(S_j)
 \quad\text{for }k-r+1\le j\le k,
 \qquad f_j=0\text{ otherwise}.
\]
For sufficiently large $k$ the active dimensions are positive;
these functions are smooth, symmetric, and compactly supported
strictly inside the required simplex.
Equation~\eqref{eq:cutoffexception} gives $w=1-O_g(k^{-1})$.
On an interior active level, normalize $f_j+Tf_{j+1}$ by
$(-1)^jP_j/\sqrt{r\alpha_j}$. The resulting expression is
\[
 \chi(S_j)-
 \sqrt{\frac{j+1}{\lambda k}}\,
       \int g(y)\chi(S_j+y/k)\,dy.
\]
Its absolute value is at most one everywhere: both terms
belong to $[0,1]$, since $j+1\le k$ and $\int g=\sqrt\lambda$.
Off the exceptional set its absolute value is at most $r/k$.
The upper and lower endpoint levels each contribute at most
$1/r$ to $v/\lambda$. At the lower endpoint this follows
from $\alpha_{j-1}/\alpha_j=j/\lambda$ and
$TP_j=\sqrt{\lambda/k}\,P_{j-1}$.
Consequently
\[
 \frac v\lambda
 \le\frac2r+(r/k)^2+O_g(k^{-1})\longrightarrow0.
\]
Since $w\to1$, fixing a sufficiently large finite $k$
proves (i).

\paragraph{Amplification.}
Now take any fixed $M>0$, retain only
$f_k=\alpha_k^{-1/2}P_k\chi(S_k)$, and set all other levels
equal to zero.
Again $w=1-O_g(k^{-1})$.
The only nonzero contributions to $v$ are the level $k$
term $\lambda w$ and the level $k-1$ deletion term.
Off an exceptional set of $P_{k-1}^2$ probability $O_g(k^{-1})$,
\[
 Tf_k=\alpha_k^{-1/2}\sqrt{M/k}\,P_{k-1};
\]
everywhere its magnitude is at most that expression.
It follows that
\[
 \lambda\alpha_{k-1}\norm{Tf_k}_2^2
   =M(1-O_g(k^{-1})),\qquad
 \frac vw\longrightarrow\lambda+M.
\]
Choose $M$ first and then a sufficiently large finite $k$
to make $v/w>H_0$. This proves (ii).
\end{proof}

\section{An empty region and a crowded region in the same class}
\label{sec:joint}

\begin{proposition}[Simultaneous suppression and occupation]
\label{prop:joint}
Let $I_A,I_B$ satisfy \eqref{eq:rowcounts}. For every fixed
integer $m\ge1$ there is $d>0$ such that, for all sufficiently
large $Q$,
\[
 \PP\{N_A=0,\ N_B\ge m\}\ge d.
\]
\end{proposition}
\begin{proof}
First choose the $B$ family using Lemma~\ref{lem:variational}
so that $v_B>(m+1)w_B>0$. Fix $G$ and the resulting finite
constant $q_B$ in Lemma~\ref{lem:tensor}.
In fact $q_B\ge v_B>0$, by \eqref{eq:majorants} and the
single-colour marked limit.
Next choose the $A$ family by suppression so that
\begin{equation}\label{eq:orderedchoice}
 \frac{v_A}{w_A}
     <\frac{v_B-mw_B}{2q_B}.
\end{equation}
This choice is legitimate because $q_B$ is independent of the
$A$ family.
Only after these two finite families have been fixed do we
let $Q$ tend to infinity.

Put $W=Z_A^2Z_B^2$ and
$E=\{N_A=0,\ N_B\ge m\}$.
Since $\ind_{N_A\ge1}\le V_A$, equations
\eqref{eq:NV} and \eqref{eq:majorants} give
\begin{equation}\label{eq:loss}
 \E[WV_B\ind_E]
 \ge \E[WV_B]-m(1+\epsilon_Q)\E W
       -(1+\epsilon_Q)\E[WV_AU_B].
\end{equation}
Indeed, on $N_B<m$ the weighted count $V_B$ is at most
$m(1+\epsilon_Q)$; the other unwanted event is bounded by
$V_B\ind_{N_A\ge1}\le(1+\epsilon_Q)V_AU_B$.
By Lemma~\ref{lem:tensor} and \eqref{eq:orderedchoice},
the lower limit of \eqref{eq:loss} is strictly larger than
\[
 \frac12w_A(v_B-mw_B)>0.
\]
H\"older's inequality yields
\[
 \E[WV_B\ind_E]
 \le(\E W^2)^{1/2}(\E V_B^4)^{1/4}\PP(E)^{1/4}.
\]
The two moments on the right are bounded for the now fixed
families. Thus $\PP(E)$ has a fixed positive lower bound for
every sufficiently large $Q$.
\end{proof}

\paragraph{The dependence issue resolved by the majorant.}
Separate positive-proportion statements about $N_A=0$ and
$N_B\ge m$ would not imply that their intersection is nonempty.
Nor would a bare Cauchy--Schwarz estimate with a fourth-moment
constant depending on the still-changing suppression family
justify the desired parameter choice.
The specific mixed moment in \eqref{eq:loss} resolves both issues:
\[
 \E[WV_AU_B]\longrightarrow v_Aq_B.
\]
The detector $U_B$ has no prime mark. One fixes its entire
$B$-constant before making $v_A/w_A$ small.
There is exactly one prime mark, in colour $A$.

\section{Exactly one prime and the physical windows}
\label{sec:singleton}

\begin{proposition}[A small singleton window]\label{prop:singleton}
Fix the detector $G$ and its constant $\kappa$ in
\eqref{eq:detectors}. Suppose the row sets satisfy
\eqref{eq:rowcounts}, with
\[
 0<\lambda_B<\frac1{2\kappa^2}.
\]
There is $d>0$ such that, for every sufficiently large $Q$,
\[
 \PP\{N_A=0,\ N_B=1\}\ge d.
\]
\end{proposition}
\begin{proof}
Choose only an $A$ family, with $w_A>0$ and
$v_A<\lambda_Bw_A/4$, by Lemma~\ref{lem:variational}.
For every nonnegative integer $n$,
$\ind_{n=1}\ge n-n(n-1)$.
Removing the event $N_A\ge1$ costs at most
$\ind_{N_A\ge1}\le V_A$. Thus, pointwise,
\[
 \ind_{\{N_A=0,N_B=1\}}
 \ge N_B-N_B(N_B-1)-V_A
 \ge N_B-R_B-V_A.
\]
Multiply by $Z_A^2$ and use
\eqref{eq:singlefirst}--\eqref{eq:singlemark}.
The lower limit of the weighted mass is at least
\[
 w_A(\lambda_B-\kappa^2\lambda_B^2)-v_A
   >\frac14w_A\lambda_B>0.
\]
Since $\E Z_A^4=O(1)$, Cauchy--Schwarz converts this positive
weighted mass into a positive ordinary proportion.
The pair term needed only an upper bound, supplied by $R_B$.
\end{proof}

\begin{proof}[Proof of Theorems~\ref{thm:cluster} and \ref{thm:single}]
Fix $0\le s<t<R$. Choose fixed
$s<s_0<t_0<t$, and put $\lambda_B=t_0-s_0$.
For the singleton theorem choose these inner endpoints close
enough that $\lambda_B<1/(2\kappa^2)$.
Set
\begin{align*}
 H_Q&=\left\lceil RL/u\right\rceil,\\
 I_B&=\left\{b\in\mathbb Z:
       \left\lceil s_0L/u\right\rceil
       \le b\le
       \left\lfloor t_0L/u\right\rfloor-1\right\},\\
 I_A&=\{1,\ldots,H_Q\}\setminus I_B .
\end{align*}
For all large $Q$, the row sets are disjoint and nonempty,
$I_B$ is one integer interval, and $I_A$ is a union of at most
two. Their normalized sizes tend to
$\lambda_B$ and $\lambda_A=R-\lambda_B>0$.
All rows lie in a fixed $O(L/u)$ range.
Every prime counted on a $B$ row belongs to
$(s_0T_Q,t_0T_Q]\subset(sT_Q,tT_Q]$.

Apply Proposition~\ref{prop:joint} for the cluster theorem or
Proposition~\ref{prop:singleton} for the singleton theorem.
These give a positive proportion with all $A$ rows empty
and the required number of prime $B$ rows.
Every prime at most $RT_Q$ in a class has the form
$a+bQ$ with $0\le b\le H_Q$.
The only row not yet controlled is $b=0$.
Discard classes whose representative $1\le a\le Q$ is itself
prime; they number at most
\[
 \pi(Q)\ll Q/L=o(\phi(Q)),
\]
since $u/L\to0$.
For each remaining class all primes below $RT_Q$ occur in
$I_B$. There are at least $m$ of them in the cluster case,
and exactly one in the singleton case.
Removing the $o(\phi(Q))$ exceptional classes preserves a
fixed positive proportion and proves both theorems.
\end{proof}

\begin{remark}[More prescribed empty windows]\label{rem:windows}
The same row-set propositions allow any fixed finite collection
of disjoint bounded empty windows and one occupied target interval.
For exact one, insert a small interval inside the target and
add its remainder to the empty region.
One can also obtain these versions directly from the main
theorems by choosing a larger finite $R$ and suppressing the
whole complement of an inner target interval below $RT_Q$.
No conclusion about two separately required occupied windows
is used.
\end{remark}


\section{Consequences, quantifiers, and remaining questions}
\label{sec:consequences}

\begin{proof}[Proof of Corollary~\ref{cor:support}]
Apply Theorem~\ref{thm:cluster} with $m=j$ and any fixed $R>t$.
There is no earlier prime, and the first $j$ primes belong to
$(sT_Q,tT_Q]$. This is the claimed positive lower mass.
\end{proof}

\begin{proof}[Proof of Corollary~\ref{cor:ratios}]
Given $\varepsilon>0$, choose
$s<s_0<t_0<t$ with $t_0/s_0<1+\varepsilon$.
Theorem~\ref{thm:cluster}, applied with these inner endpoints,
gives a positive proportion with
\[
 s_0T_Q<p_1\le p_m\le t_0T_Q,\qquad
 1\le p_m/p_1<1+\varepsilon.
\]
This proves the first limit in \eqref{eq:ratios} and its
positive-proportion refinement.
Given any $A>1$, choose $R>At$.
Theorem~\ref{thm:single} supplies a positive proportion in
$\mathcal R_Q(s,t)$ with
$p_2/p_1>R/t>A$. This proves the second limit.
The same theorem with $R>t+C$ shows that a positive
proportion in $\mathcal R_Q(s,t)$ have
$p_2-p_1>CT_Q$ for every fixed $C>0$.
\end{proof}

\begin{proof}[Proof of Corollary~\ref{cor:gp}]
Given $C>0$, apply Theorem~\ref{thm:single} with
$s=C$, $t=C+1$, and $R=C+2$.
It gives a fixed positive proportion with $p_1>CT_Q$ for
every sufficiently large $Q$. In particular the maximum
exceeds $CT_Q$ for every such $Q$.
Since $C$ is arbitrary and fixed before taking the limit,
\eqref{eq:gp} follows. The result is a limit through all
integers $Q$, not merely a limsup or a density-one assertion.
\end{proof}

\begin{corollary}[Three count regimes]\label{cor:counts}
For every fixed $c>0$, the proportions of reduced classes with
$N_Q(a,cT_Q)=0$ and with $N_Q(a,cT_Q)=1$ each have a
strictly positive lower limit. For every fixed integer $m\ge1$,
the same holds for $N_Q(a,cT_Q)\ge m$.
The constants for these assertions can be different.
\end{corollary}
\begin{proof}
The empty case follows from Corollary~\ref{cor:gp}.
For the singleton case, use Theorem~\ref{thm:single} with
any $0<s<t<c$ and $R=c$.
For the crowded case, use Theorem~\ref{thm:cluster}
with $0\le s<t<c$ and any $R>c$.
\end{proof}

\begin{corollary}[A large jump between successive class gaps]
\label{cor:gapjump}
Fix $0<s<t$ and $A>0$.
There exist an integer $M\ge2$ and $d>0$ such that, for
every sufficiently large $Q$, at least $d\phi(Q)$ reduced
classes admit an index $2\le j\le M$ for which
\[
 sT_Q<p_1(a,Q)\le p_j(a,Q)\le tT_Q,\qquad
 \frac{p_{j+1}(a,Q)-p_j(a,Q)}
      {p_j(a,Q)-p_{j-1}(a,Q)}>A .
\]
\end{corollary}
\begin{proof}
Choose $R>t+A(t-s)$ and apply Theorem~\ref{thm:cluster}
with $m=2$. Let its positive proportion be $d_0$.
For a class in that event put $j=N_Q(a,tT_Q)\ge2$.
Then $p_{j+1}>RT_Q$ and
$p_j-p_{j-1}<(t-s)T_Q$, so the displayed ratio is
greater than $(R-t)/(t-s)>A$.
By \eqref{eq:firstmean},
$\E N_Q(a,tT_Q)\le t+1$ for all sufficiently large $Q$.
Markov's inequality gives
$\PP(N_Q(a,tT_Q)>M)\le(t+1)/M$.
Choose the fixed integer $M$ so that this is less than $d_0/2$.
Removing these classes leaves the required proportion.
The bounded index $j$ may depend on the class and on $Q$.
\end{proof}

\begin{corollary}[Prescribed limiting location]\label{cor:diagonal}
For every $r>0$, one can choose reduced classes $a_Q,b_Q$
for each sufficiently large integer $Q$ such that
\begin{align}
 \frac{p_1(a_Q,Q)}{T_Q}&\longrightarrow r,&
 \frac{p_2(a_Q,Q)}{T_Q}&\longrightarrow+\infty,\label{eq:diagone}\\
 \frac{p_j(b_Q,Q)}{T_Q}&\longrightarrow r
       &&\text{for every fixed }j\ge1.\label{eq:diagmany}
\end{align}
These are qualitative existence statements, with no rate
or fixed positive proportion asserted for the selected classes.
\end{corollary}
\begin{proof}
Let $\varepsilon_k=\min(r/2,1/k)$.
For each fixed $k$, apply Theorem~\ref{thm:single} with
$s=r-\varepsilon_k$, $t=r+\varepsilon_k$, and
$R=r+k+1$.
Also apply Theorem~\ref{thm:cluster} with these endpoints
and $m=k$.
Both events are nonempty for every $Q$ beyond a finite
threshold. Choose strictly increasing integer thresholds $Q_k\ge k$ that
work for both assertions.
For $Q_k\le Q<Q_{k+1}$ choose any class from each event.
Then $k\to\infty$ with $Q$, giving \eqref{eq:diagone}.
For each fixed $j$, eventually $k\ge j$, and the first
$j$ primes of the second class all lie between
$(r-\varepsilon_k)T_Q$ and $(r+\varepsilon_k)T_Q$.
This proves \eqref{eq:diagmany}.
No uniform estimate with a growing $k$ was used.
\end{proof}

\subsection{What the proof does and does not determine}

The all-modulus conclusions have the logical form
\[
 \forall\ \text{fixed }(s,t,R,m)\quad
 \exists\,d>0,\ Q_0\quad
 \forall\,Q\ge Q_0\quad
 \#\{\text{desired reduced classes}\}\ge d\phi(Q).
\]
The constants can be extremely small and the thresholds extremely
large. No numerical percentage or explicit threshold is supplied,
and no effectivity claim is made.
The diagonal corollary only chooses thresholds after the fixed
parameter theorems are available.

There is a specific distinction between the two main theorems.
In Theorem~\ref{thm:cluster} the number of primes in the target
window may exceed $m$, so it is not valid to infer
$p_{m+1}>RT_Q$.
In Theorem~\ref{thm:single} the count is exactly one,
and therefore the assertion about $p_2$ is valid.

For exact occupancy one, the positive term in
\[
 \ind_{N=1}\ge N-N(N-1)
\]
requires only one detected prime; the two-prime term is subtracted,
so an upper bound suffices.
An attempt to force an arbitrary exact count $m\ge2$ by
the corresponding factorial-moment argument starts with an
$m$-prime lower contribution. The current proof has not
established that contribution. Similarly, neither the
singular-series average nor the two-colour sieve identities
allow us to multiply two assertions about occupied windows
as if they were independent.

Consequently this paper does not prove a Poisson point-process
limit, prescribed exact counts $m\ge2$, or arbitrary patterns
containing several separately required occupied windows.
It also does not prove the stronger Granville--Pomerance
bound $P(Q)\gg\phi(Q)(\log Q)^2$.
The primary quantitative bottleneck for that bound is
uniform control as the suppression parameter grows with $Q$.
The present transfer and functional lemmas fix all such
parameters before the arithmetic limit.

Finally, same-class prime gaps are positive multiples of $Q$.
Our small relative windows have length proportional to
$T_Q$, and $T_Q/Q=L/u\to\infty$.
The results therefore neither give bounded absolute gaps nor
assert that the relevant primes are consecutive in the ordinary
sequence of all primes.

\paragraph{Provenance and status.}
The arithmetic transfer and signed functional construction
consolidate and extend the supplied R119 manuscript
\cite{Meiji119}, with attribution to its methodological predecessor
\cite{OAI}. The same-class coupling, the singleton argument, and
the resulting localization statements are the additional arguments
presented here. This is a research manuscript, not a report of
external refereeing or formal proof verification. We make no
priority claim for the catalogue formulation of problem 971 or
for every consequence discussed above; an exhaustive priority
review is separate from checking this proof.

\appendix
\section{Support budgets and order of choices}\label{app:budgets}

The support margins used by the proof can be read directly from
the following table. Here a prime mark is an actual factor
$\vartheta(n)$; a detector is a divisor-square expression.
All entries use $\tau=\sigma=1/64$.

\begin{center}
\small
\begin{tabular}{@{}p{0.43\textwidth}ccc@{}}
\toprule
Expression & Prime marks & Total budget & Required bound\\
\midrule
$Z_i^2V_i$ & $1$ & $2\tau=1/32$ & $<1/2$\\
$Z_A^2Z_B^2V_B$ & $1$ & $4\tau=1/16$ & $<1/2$\\
$Z_A^2Z_B^2V_AU_B$ & $1$ & $4\tau+2\sigma=3/32$ & $<1/2$\\
$Z_A^2R_B$ & $0$ & $2\tau+4\sigma=3/32$ & $<1$\\
$Z_A^4Z_B^4$ & $0$ & $8\tau=1/8$ & $<1$\\
$U_B^4$ & $0$ & $8\sigma=1/8$ & $<1$\\
\bottomrule
\end{tabular}
\end{center}

For the crowded-window theorem the order is:
fix the physical windows, fix the $B$ amplification family
and $G$, evaluate its finite constant $q_B$, choose the
$A$ suppression family to satisfy \eqref{eq:orderedchoice},
and only then increase $Q$.
For the singleton theorem, fix $G$, choose the inner window
length $\lambda_B<1/(2\kappa^2)$, then suppress $A$ until
$v_A/w_A<\lambda_B/4$, and finally increase $Q$.
Thus no moment bound that depends on a changing family is
silently treated as uniform.

\begin{thebibliography}{99}

\bibitem{Meiji119}
Meiji,
\emph{Empty and crowded residue classes at the first-prime scale},
R119 working manuscript, 24 November 2026, 6 pp.
The supplied manuscript is the starting point for the present extension.

\bibitem{E49}
P. Erd\H os,
\href{https://www.renyi.hu/~p_erdos/1949-05.pdf}
{On some applications of Brun's method},
Acta Sci. Math. (Szeged) \textbf{13} (1949), 57--63.

\bibitem{E65}
P. Erd\H os,
\href{https://www.renyi.hu/~p_erdos/1965-17.pdf}
{Some recent advances and current problems in number theory},
in \emph{Lectures on Modern Mathematics III} (1965), 196--244,
especially p.~209.

\bibitem{Bloom}
T. F. Bloom,
\href{https://www.erdosproblems.com/971}{Erd\H os Problem 971},
online problem statement. Cited for the formulation and numbering,
not for a current open/solved status or priority determination.

\bibitem{GP}
A. Granville and C. Pomerance,
\href{https://math.dartmouth.edu/~carlp/least.pdf}
{On the least prime in certain arithmetic progressions},
J. London Math. Soc. (2) \textbf{41} (1990), 193--200.

\bibitem{LPS}
J. Li, K. Pratt, and G. Shakan,
\href{https://arxiv.org/abs/1607.02543}
{A lower bound for the least prime in an arithmetic progression},
Q. J. Math. \textbf{68} (2017), 729--758.

\bibitem{LeungMoments}
S.-K. Leung,
\href{https://doi.org/10.1515/forum-2024-0292}
{Moments of primes in progressions to a large modulus},
Forum Math. \textbf{38} (2026), 555--574.

\bibitem{LeungProcess}
S.-K. Leung,
\href{https://arxiv.org/abs/2407.05380}
{Pseudorandomness of primes at large scales},
Q. J. Math. \textbf{76} (2025), 251--263.

\bibitem{OAI}
OpenAI,
\href{https://github.com/openai/math/tree/main/preprints/Positive-lower-density-of-large-prime-gaps-September-25-2026}
{Positive lower density of large prime gaps},
research preprint, 25 September 2026.
The functional construction used here is reproduced in
Section~\ref{sec:variational}.

\bibitem{Bombieri}
E. Bombieri,
\emph{On the large sieve},
Mathematika \textbf{12} (1965), 201--225.

\bibitem{Vaughan}
R. C. Vaughan,
\href{https://personal.science.psu.edu/rcv4/Bombieri.pdf}
{The Bombieri--Vinogradov theorem},
expository notes, especially the maximal-endpoint formulation on p.~3.

\end{thebibliography}
\end{document}
